\uuid{aQeP}
\titre{Degré}
\auteur{bodin ; hanani ; mzari}
\organisation{exo7}
\licence{CC BY-NC-SA 4.0}
\datecreate{2020-01-20}
\niveau{L1}
\module{Algèbre}
\chapitre{Polynôme, fraction rationnelle}
\sousChapitre{Définition, degré, produit}
\difficulte{2}
\theme{Polynômes, degré}
\exo7qcmid{lille-1.polynomes.004}

\contenu{
\begin{question}
\qclassification{105.05}
\qtags{motcle=degré}
Soit $P$ un polynôme de degré $\ge 2$.
Quelles sont les assertions vraies, quel que soit le polynôme $P$ ?
\begin{answers}
    \good{$\deg( P(X) \times (X^2-X+1) ) = \deg P(X) + 2$}

    \bad{$\deg( P(X) + (X^2-X+1) ) = \deg P(X)$}

    \bad{$\deg( P(X)^2 ) = (\deg P(X))^2$}

    \good{$\deg( P(X^2) ) = 2\deg P(X)$}
\end{answers}
\begin{explanations}
On a la formule $\deg(P\times Q) = \deg P + \deg Q$ mais, il n'y a pas de formule pour la somme, car $\deg(P + Q)$ peut être strictement plus petit que $\deg P$ et $\deg Q$.
\end{explanations}
\end{question}
}
